
\section{多普勒效应}

\begin{definition}[多普勒效应]
当波源、观察者或二者相对介质运动时，观察者接收到的频率 \(\nu'\) 与波源频率 \(\nu\) 不同，这种现象称为\textbf{多普勒效应}。对机械波，频率变化必须以介质为参照。

\[
\nu'=\nu\frac{u\pm v_o}{u\mp v_s}.
\]
式中，观察者朝向波源运动取上号，远离取下号；波源朝向观察者运动时，分母取减号，远离取加号。观察者运动改变接收波峰的节奏，波源运动改变介质中的波长。
\end{definition}

\begin{exercise}[两个汽笛同时发声时怎样判出拍频]
两个汽笛 \(A,B\) 的频率都为 \(\SI{500}{Hz}\)。\(A\) 静止，\(B\) 以 \(\SI{60}{m/s}\) 向右运动。在两汽笛之间有一观察者 \(O\)，以 \(\SI{30}{m/s}\) 向右运动。已知空气中声速为 \(\SI{330}{m/s}\)。求观察者听到来自 \(A\) 的频率、来自 \(B\) 的频率以及拍频。
\end{exercise}

\begin{solution}
分开算两次多普勒，再取频差。对来自 \(A\) 的声波，\(A\) 静止而观察者远离 \(A\)，故
 \(\nu_A'=\nu\frac{u-v_o}{u} =500\times\frac{330-30}{330} \approx \SI{454.5}{Hz}.\) 
对来自 \(B\) 的声波，声源和观察者相向运动，故
 \(\nu_B'=\nu\frac{u+v_o}{u+v_s} =500\times\frac{330+30}{330+60} \approx \SI{461.5}{Hz}.\) 
拍频为
 \(f_{\text{beat}}=\abs{\nu_B'-\nu_A'} \approx \abs{461.5-454.5} \approx \SI{7.0}{Hz}.\) 
所以观察者听到来自 \(A\) 的频率约为 \(\SI{454.5}{Hz}\)，来自 \(B\) 的频率约为 \(\SI{461.5}{Hz}\)，拍频约为 \(\SI{7}{Hz}\)。
\end{solution}

\begin{exercise}[2012级期末：火车汽笛远离静止观察者的多普勒频移]
一列火车汽笛频率为 \(\SI{750}{Hz}\)，机车以时速 \(90\,\mathrm{km/h}\) 远离静止的观察者。空气中声速为 \(\SI{340}{m/s}\)，求观察者听到的声音频率。
\end{exercise}

\begin{solution}
火车速度 \(90\,\mathrm{km/h}=\SI{25}{m/s}\)，于是
 \(\nu'=\nu\frac{u}{u+v_s} =750\times\frac{340}{340+25} \approx \SI{699}{Hz}.\) 
\end{solution}

\begin{exercise}[2015-2016期末：两列相向火车中的汽笛频率]
两列火车均以 \(64.8\,\mathrm{km/h}\) 的速率迎面相向行驶。一列火车的汽笛频率为 \(\SI{600}{Hz}\)，空气中声速为 \(\SI{340}{m/s}\)。求另一列火车上的乘客听到的汽笛频率。
\end{exercise}

\begin{solution}
速度 \(64.8\,\mathrm{km/h}=\SI{18}{m/s}\)。两列车速率相同，故 \(v_o=v_s=\SI{18}{m/s}\)，于是
 \(\nu'=\nu\frac{u+v_o}{u-v_s} =600\times\frac{340+18}{340-18} =600\times\frac{358}{322} \approx\SI{667}{Hz}.\) 
\end{solution}
